\section{Routing Heads and Counting Circuits}
\label{sec:routing}

\subsection{Three routing roles}

The state analysis identifies what is represented; it does not identify how source information reaches that state. We use attention only to rank candidates on a discovery split, then assign a functional name only after query-local causal tests.

Let \(A_{qj}^{\ell h}\) be attention from semantic query \(q\) to key token \(j\) in layer \(\ell\), query head \(h\). For source span \(S_i\), define its attention mass
\begin{equation}
a_i^{\ell h}(q)=\sum_{j\in S_i}A_{qj}^{\ell h}.
\label{eq:source-mass}
\end{equation}
All source scores are averaged within an example before averaging across examples, so long traces do not dominate. Exact spans, query rows, and missing/alignment rules are frozen before head discovery.

\paragraph{Broad-source routing score.}
At the Direct pre-answer query \(q_D\), let
\[
\mathcal M(q_D)=\sum_{i=1}^{N}a_i(q_D),
\qquad
r_i(q_D)=\frac{a_i(q_D)}{\mathcal M(q_D)+\epsilon},
\]
and define normalized source entropy
\[
H_N(q_D)=
-\frac{\sum_{i=1}^{N}r_i(q_D)\log(r_i(q_D)+\epsilon)}
{\log N}.
\]
The discovery score is
\begin{equation}
B_{\ell h}=\E[\mathcal M(q_D)H_N(q_D)].
\label{eq:broad-score}
\end{equation}
\(B_{\ell h}\) is high only when a head places non-negligible total mass over multiple valid sources. We also report \(\mathcal M\), \(H_N\), source coverage, and the union coverage of top-\(k\) head bundles. \(B\) is not used for \(N<2\).

\paragraph{Matched-source routing score.}
At the CoT retrieval query \(q_k^{\rm ret}\), the intended source is \(S_{\pi_k}\). The primary ranking statistic is raw matched-source mass
\begin{equation}
T_{\ell h}^{\rm raw}
=
\E_x\!\left[
\frac{1}{K_x}\sum_{k=1}^{K_x}
a_{\pi_k}^{\ell h}(q_k^{\rm ret})
\right].
\label{eq:target-score}
\end{equation}
We additionally report conditional selectivity
\begin{equation}
T_{\ell h}^{\rm sel}
=
\E_{x,k}\!\left[
\frac{a_{\pi_k}^{\ell h}(q_k^{\rm ret})}
{\sum_i a_i^{\ell h}(q_k^{\rm ret})+\epsilon}
\right],
\label{eq:target-selectivity}
\end{equation}
source-subset top-1, and the correct-versus-best-wrong QK margin. Raw mass prevents a head with negligible source attention from ranking highly merely because its tiny mass is conditionally selective.

\paragraph{Next-source routing score.}
At item end \(q_k^{\rm end}\), for a canonical ordering with a unique next unvisited source, define
\begin{equation}
S_{\ell h}^{\rm next}
=
\E_x\!\left[
\frac{1}{K_x-1}
\sum_{k<K_x}
a_{\pi_{k+1}}^{\ell h}(q_k^{\rm end})
\right].
\label{eq:next-score}
\end{equation}
We call a high-scoring component a \emph{next-source routing candidate}, not a successor head. The stronger term \emph{continue/stop control head} is reserved for a component whose bidirectional patch changes both the next source and whether generation continues. Native traces may not have a unique gold order; canonical traces provide the primary score, while aligned free-running traces validate natural execution. The role is distinct from classic successor heads that implement an ordered-token map such as \(k\mapsto k+1\)~\citep{gould2023successor}.

\begin{table}[H]
\centering
\small
\caption{The three headline routing roles. A score nominates candidates; the final column is required for causal naming.}
\label{tab:head-roles}
\begin{tabularx}{\textwidth}{@{}P{.19\textwidth}P{.17\textwidth}P{.20\textwidth}P{.12\textwidth}Y@{}}
\toprule
Role & Query & Registered keys & Score & Causal certification \\
\midrule
Direct broad-source
  & \(q_D\)
  & all valid prompt sources
  & \(B_{\ell h}\)
  & multiple source-specific effects pass through one bundle; local ablation damages \(Z_D\) and the answer; clean patch restores both \\
CoT matched-source
  & \(q_k^{\rm ret}\)
  & intended source \(S_{\pi_k}\)
  & \(T_{\ell h}^{\rm raw}\)
  & same-position identity corruption is repaired, and counterfactual pointer changes retarget the same bundle \\
CoT next-source
  & \(q_k^{\rm end}\)
  & next unvisited source
  & \(S_{\ell h}^{\rm next}\)
  & short/long-prefix patches change next-source and continue/close decisions in both directions and alter rollout stopping \\
\bottomrule
\end{tabularx}
\end{table}

Final trace-to-answer attention is analyzed as a readout route in the appendix, not as a fourth headline mechanism. High trace mass does not imply that a head transports a scalar total; final-total transport is already tested at \(Z_C\) in \cref{sec:representations}.

\begin{figure}[H]
\centering
\fbox{\parbox[c][6.0cm][c]{0.94\textwidth}{
\centering
\textbf{TODO Figure: Routing roles and causal certification}\\[0.5em]
\begin{minipage}[t]{.30\linewidth}\centering
\textit{Broad-source at \(q_D\)}\\
Qwen: source-wide heatmap\\
Gemma: source-wide heatmap
\end{minipage}\hfill
\begin{minipage}[t]{.30\linewidth}\centering
\textit{Matched-source}\\
Qwen: step--source diagonal\\
Gemma: step--source diagonal
\end{minipage}\hfill
\begin{minipage}[t]{.30\linewidth}\centering
\textit{Next-source}\\
Qwen: item to next source\\
Gemma: item to next source
\end{minipage}\\[0.8em]
\small Under each heatmap: discovery score, query-local ablation damage, clean-patch recovery, and matched-random interval. Show source spans rather than full-token matrices, use one shared scale within each role, and mark the intervened query row.
}}
\caption{Planned attention-and-causality figure. Heatmaps describe candidate routing; the aligned ablation and restoration panels determine function. Head numbers need not match across architectures.}
\label{fig:head-todo}
\end{figure}

\subsection{One causal certification ladder}

Each candidate role uses the same evidence ladder:
\begin{equation}
\text{attention signature}
\longrightarrow
\text{query-local necessity}
\longrightarrow
\text{clean restoration}
\longrightarrow
\text{free-running effect}.
\label{eq:head-ladder}
\end{equation}
For a clean/corrupt endpoint \(m\), normalized recovery is
\begin{equation}
\operatorname{Rec}
=
\frac{m_{\rm patch}-m_{\rm corrupt}}
{m_{\rm clean}-m_{\rm corrupt}}.
\label{eq:recovery}
\end{equation}
We report the raw numerator and denominator and compute recovery only when the clean--corrupt gap exceeds a frozen minimum. Endpoints are complete source-identity margins, continue-versus-close margins, final-state metrics, or complete numeric-string probabilities; no mechanism conclusion relies on one next-token logit.

Interventions proceed from coarse to fine: residual/block output, attention output versus MLP output, and then query-local pre-output-projection head slices. This ordering matters for grouped-query attention in Qwen and Gemma: shared key/value heads do not make a numerical query-head index a portable functional unit. Pattern, source-value, pre-\(o_{\rm proj}\) head output, and post-layer residual patches are separated. Mean replacement at the registered query is primary; zero and global ablations are sensitivity analyses.

\subsection{Direct broad-source circuit}

The Direct experiment uses length-matched source edits that add, delete, or validity-flip one record while holding candidate ordinal, syntax, and depth fixed. A broad-source bundle must carry independent effects from multiple sources into \(Z_D\). A set of narrow heads with broad union coverage is allowed; the claim is about the route bundle, not the visual breadth of one head. Multi-hop routes through prompt summary tokens are also tested, because absence of direct answer-to-source attention does not imply absence of broad retrieval.

\begin{todobox}
\textbf{H1-D---Direct broad-source route.}
Split each model's 300 dense prompts into discovery and reporting sets balanced by count. On discovery data, rank single heads and cumulative bundles by \(B_{\ell h}\), total source mass, and union coverage. On reporting data, apply fixed-length source addition/deletion, correct-identity/wrong-attribute, and matched non-target edits. At \(q_D\), run query-local mean ablation and clean-to-corrupt restoration for attention pattern, source value, pre-\(o_{\rm proj}\) head output, attention output, and residual state. Primary endpoints are the number of distinct valid sources whose effects pass through the bundle, change/recovery of the frozen \(Z_D\) total-state margin, and change/recovery of the complete parsed-count distribution. Compare with same-layer score-matched random bundles, wrong query rows, wrong source spans, and candidate summary-token paths. Assign ``causal broad-source route'' only if all three endpoints replicate in both co-primary models.
\end{todobox}

\subsection{CoT matched-source circuit}

The matched-source experiment changes the identity of the intended source while preserving total count, trace step, source position, and token length. A true targeted route must follow semantic source identity rather than the visible index or absolute position. We therefore use index-free and permuted-index traces and source swaps that keep the answer unchanged.

\begin{todobox}
\textbf{H1-C---CoT matched-source route.}
Construct at least 240 same-count, same-position, same-length clean/corrupt pairs per model. Change \(\pi_k\) while preserving all other source locations. Select the top-\(K\) bundle by \(T^{\rm raw}\), conditional selectivity, and QK margin on discovery data; freeze \(K\) before reporting. At \(q_k^{\rm ret}\), ablate and patch attention pattern, matched-source value, pre-\(o_{\rm proj}\) head output, attention output, and residual state. Primary endpoints are correct-versus-best-wrong source margin, post-retrieval source-identity decoding, and generated item identity. Wrong-source donors, wrong query rows, same-layer random bundles, index permutations, and non-source tokens are controls. The role is certified only if clean patches recover identity and the same bundle retargets under a pointer/source swap. Repeat on free-running prefixes before the first trace error.
\end{todobox}

\subsection{Next-source and continuation control}

A high \(S^{\rm next}\) score may reflect attention to a likely future source without controlling it. We therefore construct nested prefixes that share the same current item but differ in whether another source remains. This creates both continue and terminal decisions at the same semantic and positional anchor.

\begin{todobox}
\textbf{H2---next-source and continue/stop control.}
Using canonical traces, construct at least 240 nested short/long prefix pairs with the same current item, visible index, token length, and item-end position but different remaining-source sets. At \(q_k^{\rm end}\), patch the discovery-ranked next-source bundle from continue to terminal examples and in the reverse direction. Measure next-source QK/source margin, generated next-item identity, continue-versus-close margin, and free-running termination location. Separately patch \(P_{C,k}\), trace-token KV states, and the candidate head output to determine whether control comes from an internal progress register or external ledger. Same-layer random heads, wrong remaining sets, wrong rows, and equal-length padding are controls. Use ``continue/stop control head'' only after both directions and both co-primary models pass; otherwise retain ``next-source routing candidate.''
\end{todobox}

\subsection{Cross-model role alignment}

Qwen3-8B is used for discovery of the analysis workflow, not for fixing Gemma head numbers. Before inspecting Gemma effects, we freeze semantic queries, source spans, ranking metrics, bundle-size selection, causal endpoints, and negative controls. Gemma candidates are then discovered independently and evaluated on a disjoint split. A cross-model replication means the same functional role at a comparable relative-depth band produces the same causal endpoint; it does not require homologous indices or identical attention pictures.
